\documentclass[12pt, ukrainian]{article}
\usepackage[a4paper]{geometry}
\setlength\parindent{0mm}
\geometry{top=1.1cm,bottom=1.1cm,left=1.1cm,right=1.1cm}
\pagestyle{empty}
\usepackage[T2A]{fontenc}
\usepackage[utf8]{inputenc}
\usepackage[ukrainian]{babel}
\usepackage{graphicx}
\usepackage{caption}
\DeclareCaptionFormat{nocaption}{}
\captionsetup{format=nocaption,aboveskip=0pt,belowskip=0pt}
\usepackage[Export]{adjustbox} 
\adjustboxset{max size={0.9\linewidth}{0.9\paperheight}}
\begin{document}
%begin
%beginitem
\begin{Large}

1. Написати функцію, що для цілого числа знаходить простий дільник, який входить у розклад числа в найбільшому степені. Якщо таких кілька, то повернути найменший. \\ Наприклад, для числа $2^{9}{\cdot}3^5{\cdot}7^9{\cdot}11$ функція має повернути $2$.\\	Якість коду також оцінюється. 



\vspace{10mm}

%enditem
\newpage

%beginitem
2. 
Написати функцію, що виводить послідовність \underline{цілих} чисел
$$\scalebox{1.2}{$(-1)^{k}C_{3k}^{k}$}, \quad k=2,\ldots,n$$
Оцінюється економність обчислень. Використовувати ** та виклики функцій заборонено.

%enditem
\newpage

%end
%end



%begin
%beginitem
1. Написати функцію, що для цілого числа знаходить простий дільник, який входить у розклад числа в найбільшому степені. Якщо таких кілька, то повернути найбільший. \\ Наприклад, для числа $2^{9}{\cdot}3^5{\cdot}7^9{\cdot}11$ функція має повернути $7$.\\	Якість коду також оцінюється. 



\vspace{10mm}

%enditem
\newpage

%beginitem
2. 
Написати функцію, що виводить послідовність \underline{цілих} чисел
$$\scalebox{1.2}{$C_{3k}^{k}$}, \quad k=2,\ldots,n$$
Оцінюється економність обчислень. Використовувати ** та виклики функцій заборонено.

%enditem
\newpage

%end
%end



%begin
%beginitem
1. Написати функцію, що для цілого числа знаходить кількість його простих дільників, що входять у його розклад не більш ніж у 5му степеню. \\ Наприклад, для числа $3^5{\cdot}7^9{\cdot}11{\cdot}23^{2}$ таких дільників два – $11$ та $23$, тому функція має повернути $2$.\\	Якість коду також оцінюється. 



\vspace{10mm}

%enditem
\newpage

%beginitem
2. 
Написати функцію, що виводить послідовність \underline{цілих} чисел
$$\scalebox{1.2}{$C_{n+2k}^{k}$}, \quad k=1,\ldots,n$$
Оцінюється економність обчислень. Використовувати ** та виклики функцій заборонено.

%enditem
\newpage

%end
%end



%begin
%beginitem
1. Написати функцію, що для інтервалу $[a, b)$ знаходить число з найбільшою кількістю різних простих дільників. Якщо їх кілька, то цікавить найбільше.\\	Якість коду також оцінюється. 



\vspace{10mm}

%enditem
\newpage

%beginitem
2. 
Написати функцію, що виводить послідовність \underline{цілих} чисел
$$\scalebox{1.2}{$C_{2n+k+1}^{n+k}$}, \quad k=0,1,\ldots,n$$
Оцінюється економність обчислень. Використовувати ** та виклики функцій заборонено.

%enditem
\newpage

%end
%end



%begin
%beginitem
1. Написати функцію, що для інтервалу $[a, b)$ знаходить число з найбільшою кількістю різних простих дільників, що містять цифру 7 у своєму десятковому запису. Якщо їх кілька, то цікавить найменше.\\	Якість коду також оцінюється. 



\vspace{10mm}

%enditem
\newpage

%beginitem
2. 
Написати функцію, що виводить послідовність \underline{цілих} чисел
$$\scalebox{1.2}{$C_{n+k}^{k+1}$}, \quad k=1,\ldots,n$$
Оцінюється економність обчислень. Використовувати ** та виклики функцій заборонено.

%enditem
\newpage

%end
%end



%begin
%beginitem
1. Написати функцію, що для цілого числа знаходить суму його простих дільників, що входять у його розклад не більш ніж у 5му степеню. \\ Наприклад, для числа $3^5{\cdot}7^9{\cdot}11{\cdot}23^{2}$ таких дільників два – $11$ та $23$, тому функція має повернути $34$.\\	Якість коду також оцінюється. 



\vspace{10mm}

%enditem
\newpage

%beginitem
2. 
Написати функцію, що виводить послідовність \underline{цілих} чисел
$$\scalebox{1.2}{$(-1)^{k}\frac{C_{2k}^k}{k+1}$}, \quad k=0,1,\ldots,n$$
Оцінюється економність обчислень. Використовувати ** та виклики функцій заборонено.

%enditem
\newpage

%end
%end



%begin
%beginitem
1. Написати функцію, що для цілого числа знаходить кількість його простих дільників, що входять у його розклад рівно у другому степеню. \\ Наприклад, для числа $3^2{\cdot}7^9{\cdot}11{\cdot}23^{2}$ таких дільників два – $3$ та $23$, тому функція має повернути $2$.\\	Якість коду також оцінюється. 



\vspace{10mm}

%enditem
\newpage

%beginitem
2. 
Написати функцію, що виводить послідовність \underline{цілих} чисел
$$\scalebox{1.2}{$C_{n+k}^{2k}$}, \quad k=0,1,\ldots,n$$
Оцінюється економність обчислень. Використовувати ** та виклики функцій заборонено.

%enditem
\newpage

%end
%end



%begin
%beginitem
1. Написати функцію, що для цілого числа знаходить кількість простих дільників, які входять у розклад числа в найбільшому степені. \\ Наприклад, для числа $3^5{\cdot}7^9{\cdot}11{\cdot}23^{9}$ таких дільників два – $23$ та $7$, тому функція має повернути $7$.\\	Якість коду також оцінюється. 



\vspace{10mm}

%enditem
\newpage

%beginitem
2. 
Написати функцію, що виводить послідовність \underline{цілих} чисел
$$\scalebox{1.2}{$(-1)^{k}C_{n+k}^{2k}$}, \quad k=1,\ldots,n$$
Оцінюється економність обчислень. Використовувати ** та виклики функцій заборонено.

%enditem
\newpage

%end
%end



%begin
%beginitem
1. Написати функцію, яка для цілого числа знаходить кількість його різних простих дільників, що мають у десятковому запису цифру 1.\\	Якість коду також оцінюється. 



\vspace{10mm}

%enditem
\newpage

%beginitem
2. 
Написати функцію, що виводить послідовність \underline{цілих} чисел
$$\scalebox{1.2}{$C_{2n+k+1}^{n+k}$}, \quad k=1,\ldots,n$$
Оцінюється економність обчислень. Використовувати ** та виклики функцій заборонено.

%enditem
\newpage

%end
%end



%begin
%beginitem
1. Написати функцію, що для інтервалу $[a, b)$ знаходить число з найбільшою кількістю різних простих дільників. Якщо їх кілька, то цікавить найменше.\\	Якість коду також оцінюється. 



\vspace{10mm}

%enditem
\newpage

%beginitem
2. 
Написати функцію, що виводить послідовність \underline{цілих} чисел
$$\scalebox{1.2}{$\frac{C_{2k}^k}{k+1}$}, \quad k=0,1,\ldots,n$$
Оцінюється економність обчислень. Використовувати ** та виклики функцій заборонено.

%enditem
\newpage

%end
%end



%begin
%beginitem
1. Написати функцію, що для інтервалу $[a, b)$ знаходить число з найбільшою кількістю різних простих дільників, що містять цифру 7 у своєму десятковому запису. Якщо їх кілька, то цікавить найбільше.\\	Якість коду також оцінюється. 



\vspace{10mm}

%enditem
\newpage

%beginitem
2. 
Написати функцію, що виводить послідовність \underline{цілих} чисел
$$\scalebox{1.2}{$(-1)^{k}C_{n+k}^{k}$}, \quad k=0,1,\ldots,n$$
Оцінюється економність обчислень. Використовувати ** та виклики функцій заборонено.

%enditem
\newpage

%end
%end



%begin
%beginitem
1. Написати функцію, яка для цілого числа знаходить суму його простих дільників.\\	Якість коду також оцінюється. 



\vspace{10mm}

%enditem
\newpage

%beginitem
2. 
Написати функцію, що виводить послідовність \underline{цілих} чисел
$$\scalebox{1.2}{$C_{2k-1}^{k}$}, \quad k=2,\ldots,n$$
Оцінюється економність обчислень. Використовувати ** та виклики функцій заборонено.

%enditem
\newpage

%end
%end



%begin
%beginitem
1. Написати функцію, що для інтервалу $[a, b)$ знаходить число з найбільшою кількістю різних простих дільників. Якщо їх кілька, то цікавить найбільше.\\	Якість коду також оцінюється. 



\vspace{10mm}

%enditem
\newpage

%beginitem
2. 
Написати функцію, що виводить послідовність \underline{цілих} чисел
$$\scalebox{1.2}{$C_{n+k}^{k}$}, \quad k=1,\ldots,n$$
Оцінюється економність обчислень. Використовувати ** та виклики функцій заборонено.

%enditem
\newpage

%end
%end



%begin
%beginitem
1. Написати функцію, що для цілого числа знаходить простий дільник, який входить у розклад числа в найбільшому степені. Якщо таких кілька, то повернути найменший. \\ Наприклад, для числа $2^{9}{\cdot}3^5{\cdot}7^9{\cdot}11$ функція має повернути $2$.\\	Якість коду також оцінюється. 



\vspace{10mm}

%enditem
\newpage

%beginitem
2. 
Написати функцію, що виводить послідовність \underline{цілих} чисел
$$\scalebox{1.2}{$(-1)^{k}C_{n+2k}^{k+1}$}, \quad k=0,1,\ldots,n$$
Оцінюється економність обчислень. Використовувати ** та виклики функцій заборонено.

%enditem
\newpage

%end
%end



%begin
%beginitem
1. Написати функцію, яка для цілого числа знаходить суму його простих дільників.\\	Якість коду також оцінюється. 



\vspace{10mm}

%enditem
\newpage

%beginitem
2. 
Написати функцію, що виводить послідовність \underline{цілих} чисел
$$\scalebox{1.2}{$(-1)^{k}C_{2n+k}^{n+k}$}, \quad k=1,\ldots,n$$
Оцінюється економність обчислень. Використовувати ** та виклики функцій заборонено.

%enditem
\newpage

%end
%end



%begin
%beginitem
1. Написати функцію, що для цілого числа знаходить простий дільник, який входить у розклад числа в найбільшому степені. Якщо таких кілька, то повернути найбільший. \\ Наприклад, для числа $2^{9}{\cdot}3^5{\cdot}7^9{\cdot}11$ функція має повернути $7$.\\	Якість коду також оцінюється. 



\vspace{10mm}

%enditem
\newpage

%beginitem
2. 
Написати функцію, що виводить послідовність \underline{цілих} чисел
$$\scalebox{1.2}{$C_{n+k}^{k+1}$}, \quad k=0,1,\ldots,n$$
Оцінюється економність обчислень. Використовувати ** та виклики функцій заборонено.

%enditem
\newpage

%end
%end



%begin
%beginitem
1. Написати функцію, що для цілого числа знаходить кількість простих дільників, які входять у розклад числа в найбільшому степені. \\ Наприклад, для числа $3^5{\cdot}7^9{\cdot}11{\cdot}23^{9}$ таких дільників два – $23$ та $7$, тому функція має повернути $7$.\\	Якість коду також оцінюється. 



\vspace{10mm}

%enditem
\newpage

%beginitem
2. 
Написати функцію, що виводить послідовність \underline{цілих} чисел
$$\scalebox{1.2}{$(-1)^{k}C_{n+2k}^{k+1}$}, \quad k=1,\ldots,n$$
Оцінюється економність обчислень. Використовувати ** та виклики функцій заборонено.

%enditem
\newpage

%end
%end



%begin
%beginitem
1. Написати функцію, що для цілого числа знаходить кількість його простих дільників, що входять у його розклад не більш ніж у 5му степеню. \\ Наприклад, для числа $3^5{\cdot}7^9{\cdot}11{\cdot}23^{2}$ таких дільників два – $11$ та $23$, тому функція має повернути $2$.\\	Якість коду також оцінюється. 



\vspace{10mm}

%enditem
\newpage

%beginitem
2. 
Написати функцію, що виводить послідовність \underline{цілих} чисел
$$\scalebox{1.2}{$(-1)^{k}C_{n+k}^{k}$}, \quad k=1,\ldots,n$$
Оцінюється економність обчислень. Використовувати ** та виклики функцій заборонено.

%enditem
\newpage

%end
%end



%begin
%beginitem
1. Написати функцію, що для цілого числа знаходить суму його простих дільників, що входять у його розклад не більш ніж у 5му степеню. \\ Наприклад, для числа $3^5{\cdot}7^9{\cdot}11{\cdot}23^{2}$ таких дільників два – $11$ та $23$, тому функція має повернути $34$.\\	Якість коду також оцінюється. 



\vspace{10mm}

%enditem
\newpage

%beginitem
2. 
Написати функцію, що виводить послідовність \underline{цілих} чисел
$$\scalebox{1.2}{$(-1)^{k}C_{n+k}^{2k-1}$}, \quad k=1,\ldots,n$$
Оцінюється економність обчислень. Використовувати ** та виклики функцій заборонено.

%enditem
\newpage

%end
%end



%begin
%beginitem
1. Написати функцію, яка для цілого числа знаходить кількість його різних простих дільників, що мають у десятковому запису цифру 1.\\	Якість коду також оцінюється. 



\vspace{10mm}

%enditem
\newpage

%beginitem
2. 
Написати функцію, що виводить послідовність \underline{цілих} чисел
$$\scalebox{1.2}{$(-1)^{k}C_{2n+k+1}^{n+k}$}, \quad k=1,\ldots,n$$
Оцінюється економність обчислень. Використовувати ** та виклики функцій заборонено.

%enditem
\newpage

%end
%end



%begin
%beginitem
1. Написати функцію, що для інтервалу $[a, b)$ знаходить число з найбільшою кількістю різних простих дільників, що містять цифру 7 у своєму десятковому запису. Якщо їх кілька, то цікавить найменше.\\	Якість коду також оцінюється. 



\vspace{10mm}

%enditem
\newpage

%beginitem
2. 
Написати функцію, що виводить послідовність \underline{цілих} чисел
$$\scalebox{1.2}{$(-1)^{k}\frac{(2k+1)!!}{(k+1)!(k-1)!}$}, \quad k=2,\ldots,n$$
Оцінюється економність обчислень. Використовувати ** та виклики функцій заборонено.

%enditem
\newpage

%end
%end



%begin
%beginitem
1. Написати функцію, що для інтервалу $[a, b)$ знаходить число з найбільшою кількістю різних простих дільників. Якщо їх кілька, то цікавить найменше.\\	Якість коду також оцінюється. 



\vspace{10mm}

%enditem
\newpage

%beginitem
2. 
Написати функцію, що виводить послідовність \underline{цілих} чисел
$$\scalebox{1.2}{$(-1)^{k}C_{3k-1}^{k}$}, \quad k=2,\ldots,n$$
Оцінюється економність обчислень. Використовувати ** та виклики функцій заборонено.

%enditem
\newpage

%end
%end



%begin
%beginitem
1. Написати функцію, що для інтервалу $[a, b)$ знаходить число з найбільшою кількістю різних простих дільників, що містять цифру 7 у своєму десятковому запису. Якщо їх кілька, то цікавить найбільше.\\	Якість коду також оцінюється. 



\vspace{10mm}

%enditem
\newpage

%beginitem
2. 
Написати функцію, що виводить послідовність \underline{цілих} чисел
$$\scalebox{1.2}{$(-1)^{k}C_{2n+k}^{n+k}$}, \quad k=0,1,\ldots,n$$
Оцінюється економність обчислень. Використовувати ** та виклики функцій заборонено.

%enditem
\newpage

%end
%end



%begin
%beginitem
1. Написати функцію, що для цілого числа знаходить кількість його простих дільників, що входять у його розклад рівно у другому степеню. \\ Наприклад, для числа $3^2{\cdot}7^9{\cdot}11{\cdot}23^{2}$ таких дільників два – $3$ та $23$, тому функція має повернути $2$.\\	Якість коду також оцінюється. 



\vspace{10mm}

%enditem
\newpage

%beginitem
2. 
Написати функцію, що виводить послідовність \underline{цілих} чисел
$$\scalebox{1.2}{$C_{2k+1}^{k}$}, \quad k=2,\ldots,n$$
Оцінюється економність обчислень. Використовувати ** та виклики функцій заборонено.

%enditem
\newpage

%end
%end



%begin
%beginitem
1. Написати функцію, що для цілого числа знаходить простий дільник, який входить у розклад числа в найбільшому степені. Якщо таких кілька, то повернути найбільший. \\ Наприклад, для числа $2^{9}{\cdot}3^5{\cdot}7^9{\cdot}11$ функція має повернути $7$.\\	Якість коду також оцінюється. 



\vspace{10mm}

%enditem
\newpage

%beginitem
2. 
Написати функцію, що виводить послідовність \underline{цілих} чисел
$$\scalebox{1.2}{$\frac{(2k)!!}{(k+1)!(k-1)!}$}, \quad k=2,\ldots,n$$
Оцінюється економність обчислень. Використовувати ** та виклики функцій заборонено.

%enditem
\newpage

%end
%end



%begin
%beginitem
1. Написати функцію, що для цілого числа знаходить кількість простих дільників, які входять у розклад числа в найбільшому степені. \\ Наприклад, для числа $3^5{\cdot}7^9{\cdot}11{\cdot}23^{9}$ таких дільників два – $23$ та $7$, тому функція має повернути $7$.\\	Якість коду також оцінюється. 



\vspace{10mm}

%enditem
\newpage

%beginitem
2. 
Написати функцію, що виводить послідовність \underline{цілих} чисел
$$\scalebox{1.2}{$C_{n+k}^{k}$}, \quad k=0,1,\ldots,n$$
Оцінюється економність обчислень. Використовувати ** та виклики функцій заборонено.

%enditem
\newpage

%end
%end



%begin
%beginitem
1. Написати функцію, що для цілого числа знаходить простий дільник, який входить у розклад числа в найбільшому степені. Якщо таких кілька, то повернути найменший. \\ Наприклад, для числа $2^{9}{\cdot}3^5{\cdot}7^9{\cdot}11$ функція має повернути $2$.\\	Якість коду також оцінюється. 



\vspace{10mm}

%enditem
\newpage

%beginitem
2. 
Написати функцію, що виводить послідовність \underline{цілих} чисел
$$\scalebox{1.2}{$(-1)^{k}C_{n+2k}^{k}$}, \quad k=1,\ldots,n$$
Оцінюється економність обчислень. Використовувати ** та виклики функцій заборонено.

%enditem
\newpage

%end
%end



%begin
%beginitem
1. Написати функцію, що для інтервалу $[a, b)$ знаходить число з найбільшою кількістю різних простих дільників. Якщо їх кілька, то цікавить найменше.\\	Якість коду також оцінюється. 



\vspace{10mm}

%enditem
\newpage

%beginitem
2. 
Написати функцію, що виводить послідовність \underline{цілих} чисел
$$\scalebox{1.2}{$\frac{(2k)!!}{(k!)^2}$}, \quad k=1,\ldots,n$$
Оцінюється економність обчислень. Використовувати ** та виклики функцій заборонено.

\end{Large}

%end
%end



\end{document}

